\section{Scriptural Date and Location}
\label{sec:date-location}
{\small Each passage the Mass takes from Scripture, once, in canonical order,
with its setting, attribution and dating evidence.\par}
\medskip

\begin{dossiertable}
Introit (verse) & Ps 77:1 (Heb.\ 78) & Asaph, by the title, to Israel as ``my
people''; a history from Egypt to Sion and David (vv.~68--70) &
\chronodate{introit}{\chronologyannotation{introit}}\\
\cmidrule(lr){1-4}
\dossierprose{Titled in the Vulgate \latin{Intellectus Asaph}, ``Understanding
for Asaph'', which the Missal's verse leaves out; the title gives no occasion or
date. The displayed date is the bound that the \work{New American Bible,
Revised Edition} sets for the whole Psalter: its introduction to the Psalms
holds that no psalm is as late as the Maccabean period and that none can be
dated securely. The antiphon, printed without a citation and standing in no
verse of the Vulgate, has no dossier of its own.}
\midrule
Alleluia & Ps 104:1 (Heb.\ 105) & Israel, bidden to declare God's deeds among
the Gentiles; sung before the ark at Jerusalem (1 Par 15:3; 16:7--8) &
\chronodate{alleluia}{\chronologyannotation{alleluia}}\\
\cmidrule(lr){1-4}
\dossierprose{Titled \latin{Alleluia} in the Vulgate, which the Missal's verse
leaves out. Paralipomenon sets these words in the praise Asaph sang on the day
the ark was set in David's tent (1 Par 16:1, 7); the displayed date is the
Psalter's bound, not that day's.}
\midrule
Communion & Ps 118:4--5 (Heb.\ 119) & The psalmist, ``a sojourner on the
earth'' (v.~19), singing the law ``in the place of my pilgrimage'' (v.~54) &
\chronodate{communion}{\chronologyannotation{communion}}\\
\cmidrule(lr){1-4}
\dossierprose{Titled \latin{Alleluia}, with no author named; the Communion's
verses stand in its first alphabetic strophe, \latin{Aleph}. The displayed date
is the Psalter's bound.}
\midrule
Offertory & Ps 137:7 (Heb.\ 138) & David, by the title; praise ``in the sight
of the angels'' and worship ``towards thy holy temple'' (vv.~1--2) &
\chronodate{offertory}{\chronologyannotation{offertory}}\\
\cmidrule(lr){1-4}
\dossierprose{Titled \latin{Ipsi David}, ``For David himself'', without an
occasion; David's reign is not given as the psalm's date, and the displayed
date is the Psalter's bound.}
\midrule
Gradual & Ps 140:2 (Heb.\ 141) & David, by the title; an evening prayer beside
Israel's incense and evening sacrifice (v.~2) &
\chronodate{gradual}{\chronologyannotation{gradual}}\\
\cmidrule(lr){1-4}
\dossierprose{Titled \latin{Psalmus David}, ``A psalm of David'', again without
an occasion; David's reign is not given as its date, and the displayed date is
the Psalter's bound.}
\midrule
Gospel & Mt 22:1--14 & Writing: as Matthew left Palestine (Eusebius). Event:
Jerusalem, in the temple (21:23) & \chronodate{gospel}{\chronologyannotation{gospel}}\\
\cmidrule(lr){1-4}
\dossierevent{Narrated event: the third parable Jesus speaks in the temple, on
the day after his entry into Jerusalem (21:1--18), when the chief priests and
ancients have questioned his authority (21:23), to the chief priests and
Pharisees (21:45); St Matthew gives it no date beyond that day.}
\cmidrule(lr){1-4}
\dossierprose{Traditional attribution: St Matthew the Apostle, whom the Missal's
heading names (\latin{secundum Matthaeum}). The \work{Catholic Encyclopedia}'s
``Gospel of St.\ Matthew'' (vol.~10, 1911) reports Eusebius's tradition that
Matthew wrote his Gospel in Hebrew when he left Palestine. Its five dates,
displayed first, reckon from the Ascension, the Apostles' dispersal and St
Irenaeus, or give the Catholic critics' preference, and it settles on none, the
ancient writers being ``at variance''. The last date, c.~A.D.~50, is Durand's in
``The New Testament'' (vol.~14, 1912) for the lost Aramaic original, of which
he adds, ``We know nothing definite of the date of its being rendered into
Greek.''}
\midrule
Epistle & Eph 4:23--28 & Paul, ``a prisoner in the Lord'' (4:1), in Caesarea or
Rome; to ``the saints who are at Ephesus'' (1:1) &
\chronodate{epistle}{\chronologyannotation{epistle}}\\
\cmidrule(lr){1-4}
\dossierprose{Traditional attribution: St Paul, as the Missal's heading names
him, who names himself (1:1) and writes in chains (3:1; 6:20). The range is
Ladeuze's in the \work{Catholic Encyclopedia}'s ``Epistle to the Ephesians''
(vol.~5, 1909), where of the letters of the captivity ``whether they were
produced in Caesarea or in Rome \dots\ is still a much mooted question''. The
year is from the table in Prat's ``St.\ Paul'' (vol.~11, 1911), which sets the
letter in the Roman captivity with Philemon, Colossians and Philippians; the
articles date that captivity differently, and neither answer is settled.}
\end{dossiertable}
